% from ega4-18-fr.tex
\oldpage[IV-4]{164}

D'autre part, en vertu de (18.12.13)\footnote{Le lecteur vérifiera que (18.10.16) n'est pas utilisé dans la démonstration de (18.12.13). Si $f$ est supposé localement de présentation finie (resp. quasi-projectif), on peut remplacer (18.12.13) par (8.12.8) (resp. (8.12.6)).}, on peut (après avoir remplacé $Y$ par un voisinage ouvert affine $V$ de $y$ et $f$ par sa restriction $f^{-1}(V)\to V$) factoriser $f$ en $X\xrightarrow{j}X'\xrightarrow{g}Y$, où $g$ est fini et $j$ une immersion ouverte. Puisque $X$ est réduit, on peut remplacer $X'$ par le sous-préschéma fermé réduit d'espace sous-jacent $\overline{j(X)}$ (I, 5.2.2), donc supposer $X'$ réduit et tel que toute composante irréductible de $X'$ domine $Y$. En outre, $j(X)$ est un ouvert dense dans $X'$ et n'a qu'un nombre fini de composantes irréductibles en vertu de l'hypothèse; donc ($0_{\mathrm I}$, 1.2.7) les composantes irréductibles de $X'$ sont les adhérences de celles de $j(X)$, et par suite la fibre $g^{-1}(\eta)$ s'identifie à $f^{-1}(\eta)$. Si on suppose le théorème démontré lorsqu'en outre $f$ est un morphisme fini, on pourra l'appliquer à $g$, et comme $f^{-1}(y)$ s'identifie à un sous-préschéma de $g^{-1}(y)$, la relation (i) pour $g$ entraîne la même relation pour $f$. Supposons d'autre part que $n(y)=n$; alors, en vertu de ce qui précède, on peut appliquer la conclusion de (ii) au morphisme $g$, et en remplaçant au besoin $Y$ par un voisinage ouvert de $y$, on peut supposer $g$ étale, et il en est par suite de même de $f$; en outre, on a alors $f^{-1}(y)=g^{-1}(y)$. Comme le morphisme $g$ est fermé et que $j(X)$ est ouvert dans $X'$, il existe un voisinage $U$ de $y$ dans $Y$ tel que $j(X)\cap g^{-1}(U)=j(f^{-1}(U))$, ce qui prouve que (ii) est vrai aussi pour $f$.

On est ainsi ramené au cas où $f$ est en outre un morphisme fini. Il est immédiat que, pour prouver (i), on peut se borner au cas où $Y=\Spec(\sh{O}_{Y,y})$. Montrons qu'il en est de même pour prouver (ii). En effet, supposons prouvé que la condition $n(y)=n$ entraîne que $X\times_Y\Spec(\sh{O}_{Y,y})$ est étale et fini sur $\Spec(\sh{O}_{Y,y})$; comme cela entraîne que $f$ est plat et formellement non ramifié aux points de $f^{-1}(y)$, et que $Y$ est intègre, on conclut de (18.4.13) que $f$ est étale aux points de $f^{-1}(y)$, donc aussi dans un voisinage ouvert $V$ de $f^{-1}(y)$ dans $X$; mais puisque $f$ est un morphisme fini, donc fermé, $V$ contient un ouvert de la forme $f^{-1}(U)$, où $U$ est un voisinage ouvert de $y$ dans $Y$.

\emph{C) Réduction au cas où $Y=\Spec(\sh{O}_{Y,y})$ et $\sh{O}_{Y,y}$ est strictement local.} --- Posons $A=\sh{O}_{Y,y}$, et soit $A'={}^{hs}\!A$ le hensélisé strict de $A$ (18.8.7), qui est un anneau local intègre et géométriquement unibranche (18.8.16). Soit $Y'=\Spec(A')$, et soient $y'$ le point fermé et $\eta'$ le point générique de $Y'$; $y'$ est au-dessus de $y$ et puisque le morphisme $g:Y'\to Y$ est plat (18.8.8), $\eta'$ est au-dessus de $\eta$ (2.3.4). Posons $X'=X\times_Y Y'$, $f'=f_{(Y')}:X'\to Y'$; $f'$ est fini et puisque $g$ est plat, toute composante irréductible de $X'$ domine $Y'$ (2.3.7).

Si $n'(y')$, $n'$ et $n'_s$ sont les nombres définis pour $f'$ de la même façon que $n(y)$, $n$ et $n_s$ pour $f$, on a $n'_s=n_s$ et $n'(y')=n(y)$ (I, 6.4.8), et il revient donc au même de prouver l'assertion (i) de l'énoncé pour $f$ ou pour $f'$. En outre, si $X$ est réduit et $y$ un point normal de $Y$ (donc $A$ intégralement clos), alors $A'$ est intégralement clos (18.8.12). D'autre part il résulte de la définition (18.8.7), de la remarque (8.1.2, $a$)) et du théorème de la double limite inductive, que $A'$ est limite inductive d'anneaux $A_\alpha$ tels que les morphismes $\Spec(A_\alpha)\to\Spec(A)=Y$ soient étales; par suite $X'$ est limite projective des $X_\alpha=X\otimes_A A_\alpha$, qui sont étales sur $X$, donc réduits puisque $X$ l'est (17.5.7); par passage à la limite,

